Las curvas normalizadas confirman los dos modelos de costo: la fuerza bruta se
estabiliza en $2{,}6$ nanosegundos por nodo del árbol de estados, cuyo tamaño
es $\binom{N+M}{N}$ y no la cantidad de caminos, y la programación dinámica en
$8{,}5$ nanosegundos por estado sobre $N \cdot M \cdot V$ estados, con una
memoria que se duplica exactamente al duplicarse el área. Esa diferencia de
paradigma, y no de implementación, es la que produce una razón superior a
$2 \times 10^5$ en el mayor tamaño donde ambas corren. En cuanto al greedy, su
deterioro no es una degradación de la regla sino una acumulación: normalizada
por la cantidad de decisiones, su pérdida es constante en $2{,}00$ puntos por
movimiento en D10 y $197$ en D1000. Lo que sí depende del terreno es su tasa de
falla, un $13{,}2$ por ciento global que se concentra por completo en las
grillas con murallas, porque solo una muralla puede dejarlo sin candidatas.
